% Apollo原始多独立ST约束（子图 a）
\begin{tikzpicture}[scale=0.75, transform shape]
  \draw[-Stealth, thick] (0,0) -- (8.5,0) node[right, font=\footnotesize] {$t$};
  \draw[-Stealth, thick] (0,0) -- (0,7.5) node[above, font=\footnotesize] {$s$};

  % 三个独立 ST boundary 禁行区
  \fill[dangerred!30, draw=dangerred, thick] (0,1.0) rectangle (7.5,2.0);
  \fill[lightorange!30, draw=lightorange!80!black, thick] (1,3.5) -- (6,4.5) -- (6,5.2) -- (1,4.2) -- cycle;
  \fill[deepblue!20, draw=deepblue, thick] (2,5.8) rectangle (7.5,6.6);

  \node[dangerred, font=\footnotesize] at (-0.8, 1.5) {$B_1$};
  \node[lightorange!80!black, font=\footnotesize] at (-0.8, 3.8) {$B_2$};
  \node[deepblue, font=\footnotesize] at (-0.8, 6.2) {$B_3$};

  % 每个 boundary 的 YIELD/OVERTAKE 离散选择标注
  \node[dangerred, font=\tiny, align=center] at (3.75, 1.5) {YIELD 或 OVERTAKE?};
  \node[lightorange!80!black, font=\tiny, align=center] at (3.5, 4.85) {YIELD 或 OVERTAKE?};
  \node[deepblue, font=\tiny, align=center] at (4.75, 6.2) {YIELD 或 OVERTAKE?};

  % 底部标注
  \node[draw=dangerred, fill=dangerred!10, rounded corners, font=\footnotesize, align=center] at (4, -1.2) {$2^n$ 种 YIELD/OVERTAKE 组合\\非凸离散选择 \;\; DP 启发式搜索};

  % DP搜索路径（锯齿穿行）
  \draw[thick, deepblue, dashed, -{Stealth}]
    plot[smooth, tension=0.4] coordinates {(0,0.3) (1,0.8) (2,2.5) (3,3.0) (4,3.2) (5,5.5) (6,5.6) (7,5.7)};
  \node[deepblue, font=\tiny\bfseries] at (6.5, 5.2) {DP路径};
\end{tikzpicture}
